% The schedule: one release, four tasks, and the deadline where it was written.
\begin{tikzpicture}[font=\sffamily\scriptsize]

\draw[taxis] (0mm,-5mm) -- (122mm,-5mm) node[right, font=\tiny] {time};

\lane{40}{interrupts}
\slice{40}{0}{5}{taskisr}{timer}
\slice{40}{34}{5}{taskisr}{watermark}

\lane{32}{supervisor}
\slice{32}{0}{100}{taskblock}{blocked on its period}
\slice{32}{100}{6}{taskrun}{feed}

\lane{24}{sensor}
\slice{24}{0}{6}{taskblock}{}
\slice{24}{6}{12}{taskrun}{read FIFO}
\slice{24}{18}{21}{taskblock}{blocked on notification}
\slice{24}{39}{11}{taskrun}{read FIFO}
\slice{24}{50}{56}{taskblock}{blocked on notification}

\lane{16}{feature}
\slice{16}{0}{18}{taskblock}{}
\slice{16}{18}{16}{taskrun}{biquad}
\slice{16}{34}{16}{taskready}{ready, preempted}
\slice{16}{50}{14}{taskrun}{statistics}
\slice{16}{64}{42}{taskblock}{blocked on the stream buffer}

\lane{8}{link}
\slice{8}{0}{64}{taskblock}{blocked on the queue}
\slice{8}{64}{22}{taskrun}{AT command}
\slice{8}{86}{20}{taskblock}{awaiting the reply}

\lane{0}{idle}
\slice{0}{86}{14}{taskrun}{idle hook counts}

% the measured interval and the deadline
\draw[release] (0mm,53mm) -- (0mm,46mm);
\node[font=\tiny, anchor=west] at (1mm,55mm) {release};
\draw[deadline] (14mm,-5mm) -- (14mm,53mm);
\node[font=\tiny, text=red!70!black, anchor=west] at (16mm,55mm)
  {the deadline, written down before the first run};
\draw[tspan] (0mm,48mm) -- (6mm,48mm);
\node[tlabel, anchor=west] at (8mm,50mm) {measured: release to task running};

% legend
\slice{-16}{0}{18}{taskrun}{running}
\slice{-16}{20}{18}{taskready}{ready}
\slice{-16}{40}{18}{taskblock}{blocked}
\slice{-16}{60}{18}{taskisr}{interrupt}

\node[umlnote, text width=54mm, anchor=north west] at (0mm,-24mm)
  {The second release preempts the feature task, which becomes ready rather
   than running and resumes where it stopped. That preemption is the whole
   reason the sensor task has the higher priority.};
\node[umlnote, text width=54mm, anchor=north west] at (60mm,-24mm)
  {The shaded interval at the far left is what the million-event distribution
   measures. Everything to its right is the application doing its work, and
   none of it is allowed to move that interval.};
\end{tikzpicture}
